%! Trigonometry and Polar Coordinates — Unit 7, Lesson 7.6 — Homework
\documentclass[10pt]{article}
\usepackage{precalculus-article}
\usepackage{precalculus-boxes}
\newcommand{\TallMath}[1]{$\displaystyle #1\rule[-1.4em]{0pt}{3.2em}$}

\begin{document}

\pageheader{Unit 7, Lesson 7.6}{Homework}
\namedateperiod

\vspace{0.10in}

\begin{notesbox}{Problems 1--7}

\textbf{1.}\quad Convert each polar coordinate pair to rectangular coordinates.
Give exact values.

\vspace{4pt}
(a)\quad $(6,\, \pi/6)$:

\vspace{4pt}
$x = 6\cos(\pi/6) = $ \underline{\hspace{2.5cm}} \qquad
$y = 6\sin(\pi/6) = $ \underline{\hspace{2.5cm}} \qquad
Rectangular: $(\underline{\hspace{1.5cm}},\; \underline{\hspace{1.5cm}})$

\vspace{8pt}
(b)\quad $(4,\, 5\pi/4)$:

\vspace{4pt}
$x = $ \underline{\hspace{3cm}} \qquad $y = $ \underline{\hspace{3cm}} \qquad
Rectangular: $(\underline{\hspace{1.5cm}},\; \underline{\hspace{1.5cm}})$

\vspace{0.14in}
\textbf{2.}\quad Convert each rectangular point to polar form with $r > 0$
and $0 \le \theta < 2\pi$.

\vspace{4pt}
(a)\quad $(\sqrt{3},\, 1)$:

\vspace{4pt}
$r = $ \underline{\hspace{2.5cm}} \qquad
$\theta = $ \underline{\hspace{3cm}} (quadrant: \underline{\hspace{1cm}}) \qquad
Polar: $(\underline{\hspace{1.5cm}},\; \underline{\hspace{1.5cm}})$

\vspace{8pt}
(b)\quad $(-2,\, 2\sqrt{3})$:

\vspace{4pt}
$r = $ \underline{\hspace{2.5cm}} \qquad
$\theta = $ \underline{\hspace{3.5cm}} (quadrant: \underline{\hspace{1cm}}) \qquad
Polar: $(\underline{\hspace{1.5cm}},\; \underline{\hspace{1.5cm}})$

\vspace{0.14in}
\textbf{3.}\quad List three different polar representations of the point whose
rectangular coordinates are $(0,\, -5)$.

\vspace{4pt}
Representation 1: \underline{\hspace{4cm}} \quad
Representation 2: \underline{\hspace{4cm}} \quad
Representation 3: \underline{\hspace{4cm}}

\vspace{0.14in}
\textbf{4.}\quad Write each complex number in polar form $(r\cos\theta)+i(r\sin\theta)$.

\vspace{4pt}
(a)\quad $-3 + 3i$:

\vspace{4pt}
$r = $ \underline{\hspace{2cm}} \qquad $\theta = $ \underline{\hspace{3cm}} \qquad
Polar form: \underline{\hspace{5cm}}

\vspace{8pt}
(b)\quad $4i$:

\vspace{4pt}
$r = $ \underline{\hspace{2cm}} \qquad $\theta = $ \underline{\hspace{3cm}} \qquad
Polar form: \underline{\hspace{5cm}}

\vspace{0.14in}
\textbf{5.}\quad A polar point is given as $(2,\, \pi/4)$.

\vspace{4pt}
(a)\quad Find the rectangular coordinates.

\vspace{4pt}
Rectangular: $(\underline{\hspace{2cm}},\; \underline{\hspace{2cm}})$

\vspace{6pt}
(b)\quad Write two other polar representations of this same point.

\vspace{4pt}
Rep.\ 1: \underline{\hspace{4cm}} \qquad Rep.\ 2: \underline{\hspace{4cm}}

\vspace{0.14in}
\textbf{6. (AP-Style MC)}\quad Which of the following is the polar form of the
rectangular point $(-\sqrt{3},\, 1)$ with $r > 0$ and $0 \le \theta < 2\pi$?

\begin{enumerate}[label=(\Alph*), leftmargin=2em, itemsep=3pt]
  \item $(2,\; \pi/6)$
  \item $(2,\; 5\pi/6)$
  \item $(2,\; 7\pi/6)$
  \item $(2,\; 11\pi/6)$
  \item $(1,\; 5\pi/6)$
\end{enumerate}

Answer: \underline{\hspace{1.5cm}}

\vspace{0.12in}
\textbf{7. (AP-Style MC)}\quad The polar point $\bigl(3,\, \tfrac{2\pi}{3}\bigr)$
is the same point as which of the following?

\begin{enumerate}[label=(\Alph*), leftmargin=2em, itemsep=3pt]
  \item $\bigl(-3,\; \dfrac{2\pi}{3}\bigr)$
  \item $\bigl(3,\; \dfrac{5\pi}{3}\bigr)$
  \item $\bigl(-3,\; \dfrac{5\pi}{3}\bigr)$
  \item $\bigl(3,\; -\dfrac{\pi}{3}\bigr)$
  \item $\bigl(-3,\; \dfrac{\pi}{3}\bigr)$
\end{enumerate}

Answer: \underline{\hspace{1.5cm}}

\end{notesbox}

\vspace{0.08in}

\begin{tcolorbox}[colback=navylight, colframe=navy, arc=2mm,
    left=3mm, right=3mm, top=2mm, bottom=2mm]
\small\textbf{Preview --- Lesson 3.14: Polar Function Graphs}\\
You can now locate any point using $(r,\theta)$. In Lesson 3.14 we let $r$
be a \emph{function} of $\theta$ and graph equations like $r = 1+\cos\theta$.
The resulting curves --- cardioids, limaçons, rose curves --- are beautiful
and nearly impossible to express in rectangular form. Come ready to evaluate
a polar function at multiple angles and plot the results.
\end{tcolorbox}

\end{document}
